\begin{tomtat}
	\begin{minipage}[t]{0.3\textwidth}
		\vspace{0pt}
		\begin{tomtat}
			\begin{center}
				\textbf{ĐỀ 52}
			\end{center}
			
		\end{tomtat}
	\end{minipage}
	\hfill
	\begin{minipage}[t]{0.5\textwidth}
		\vspace{0pt}
		\begin{tomtat}
			\begin{center}
				\includegraphics[scale=0.08]{images/52}
			\end{center}
		\end{tomtat}
	\end{minipage}
\end{tomtat}
\textbf{I. PHẦN TRẮC NGHIỆM}
\begin{ex} % Bạn có thể thêm điểm ở đây, ví dụ: [(0.25 điểm)]
	Kết quả rút gọn biểu thức $A = \sqrt{12} - \sqrt{3}$ là:
	\choice{$9$}{$3$}{$\sqrt{3}$}{$2\sqrt{3}$}
	\loigiai{
		\textbf{Đáp án đúng là C.} \\
		\begin{align*}
			A &= \sqrt{12} - \sqrt{3} \\
			&= \sqrt{4 \cdot 3} - \sqrt{3} \\
			&= 2\sqrt{3} - \sqrt{3} \\
			&= \sqrt{3}
		\end{align*}
	}
\end{ex}

\begin{ex}
	Nghiệm của phương trình $2x + 1 = 11$ là:
	\choice{$5$}{$-5$}{$6$}{$-6$}
	\loigiai{
		\textbf{Đáp án đúng là A.} \\
		\begin{align*}
			2x + 1 &= 11 \\
			\Leftrightarrow 2x &= 11 - 1 \\
			\Leftrightarrow 2x &= 10 \\
			\Leftrightarrow x &= 5
		\end{align*}
	}
\end{ex}

\begin{ex}
	Đồ thị của hàm số $y = 2x^2$ đi qua điểm nào sau đây:
	\choice{$(1; 2)$}{$(2; 1)$}{$(2; 2)$}{$(2; 4)$}
	\loigiai{
		\textbf{Đáp án đúng là A.} \\
		Xét điểm A $(1; 2)$: Thay $x=1, y=2$ vào hàm số, ta có $2 = 2(1)^2 \Leftrightarrow 2 = 2$ (đúng). \\
		Vậy đồ thị hàm số đi qua điểm $(1; 2)$.
		% Lưu ý: Phần giải thích cho các điểm khác có thể không cần thiết khi đã tìm ra đáp án đúng.
		% Nếu muốn chi tiết hơn, có thể thêm lại như sau:
		% Xét điểm B $(2; 1)$: Thay $x=2, y=1$ vào hàm số, ta có $1 = 2(2)^2 \Leftrightarrow 1 = 8$ (sai). \\
		% Xét điểm C $(2; 2)$: Thay $x=2, y=2$ vào hàm số, ta có $2 = 2(2)^2 \Leftrightarrow 2 = 8$ (sai). \\
		% Xét điểm D $(2; 4)$: Thay $x=2, y=4$ vào hàm số, ta có $4 = 2(2)^2 \Leftrightarrow 4 = 8$ (sai).
	}
\end{ex}

\begin{ex}
	Gieo một con xúc xắc cân đối và đồng chất. Gọi biến cố A: ``số chấm xuất hiện trên con xúc xắc bé hơn 3''. Số kết quả thuận lợi cho biến cố A là:
	\choice{$6$}{$2$}{$3$}{$4$}
	\loigiai{
		\textbf{Đáp án đúng là B.} \\
		Các kết quả có thể xảy ra khi gieo xúc xắc là $\{1, 2, 3, 4, 5, 6\}$. \\
		Số chấm xuất hiện bé hơn 3 bao gồm các kết quả: $\{1, 2\}$. \\
		Vậy có 2 kết quả thuận lợi cho biến cố A.
	}
\end{ex}
\begin{ex}
	Cho tam giác ABC vuông tại A. Biết $AB = 6 \text{ cm}$, $AC = 10 \text{ cm}$, $\tan C$ bằng:
	\choice{$\dfrac{3}{5}$}{$\dfrac{4}{3}$}{$\dfrac{3}{4}$}{$\dfrac{5}{4}$}
	\loigiai{
		\textbf{Đáp án đúng là A.} \\
		Trong tam giác ABC vuông tại A, ta có:
		\begin{align*}
			\tan C &= \frac{\text{cạnh đối}}{\text{cạnh kề}} \\
			&= \frac{AB}{AC} \\
			&= \frac{6}{10} \\
			&= \frac{3}{5}
		\end{align*}
	}
\end{ex}

\begin{ex}
	Thể tích của một hộp sữa hình trụ có bán kính đáy $6 \text{ cm}$ và chiều cao $10 \text{ cm}$ là:
	\choice{$360\pi \text{ cm}^3$}{$600\pi \text{ cm}^3$}{$720\pi \text{ cm}^3$}{$1200\pi \text{ cm}^3$}
	\loigiai{
		\textbf{Đáp án đúng là A.} \\
		Thể tích khối trụ được tính bằng công thức $V = \pi r^2 h$, trong đó $r$ là bán kính đáy và $h$ là chiều cao.
		\begin{align*}
			V &= \pi \times (6 \text{ cm})^2 \times 10 \text{ cm} \\
			&= \pi \times 36 \text{ cm}^2 \times 10 \text{ cm} \\
			&= 360\pi \text{ cm}^3
		\end{align*}
	}
\end{ex}

\begin{ex}
	Nghiệm của bất phương trình $2x - 4 \ge 0$ là:
	\choice{$x < 2$}{$x \le 2$}{$x \ge 2$}{$x \ge -2$}
	\loigiai{
		\textbf{Đáp án đúng là C.} \\
		\begin{align*}
			2x - 4 &\ge 0 \\
			\Leftrightarrow 2x &\ge 4 \\
			\Leftrightarrow x &\ge \frac{4}{2} \\
			\Leftrightarrow x &\ge 2
		\end{align*}
	}
\end{ex}

\begin{ex}
	Cô Trang ghi lại điểm kiểm tra, đánh giá định kỳ môn Toán của một nhóm học sinh lớp 9 như sau:
	\begin{center}
		\begin{tabular}{|c|c|c|c|c|}
			\hline
			6 & 7 & 5 & 9 & 10 \\
			\hline
			5 & 8 & 10 & 5 & 10 \\
			\hline
			7 & 8 & 6 & 6 & 7 \\
			\hline
			9 & 9 & 7 & 5 & 8 \\
			\hline
		\end{tabular}
	\end{center}
	Tần số của giá trị 10 điểm là:
	\choice{$5$}{$4$}{$2$}{$3$}
	\loigiai{
		\textbf{Đáp án đúng là D.} \\
		Quan sát bảng số liệu, ta đếm số lần xuất hiện của giá trị 10 điểm:
		\begin{itemize}
			\item Dòng 1: có 1 số 10.
			\item Dòng 2: có 2 số 10.
			\item Dòng 3: không có số 10.
			\item Dòng 4: không có số 10.
		\end{itemize}
		Vậy, tần số của giá trị 10 điểm là $1 + 2 = 3$.
	}
\end{ex}


\newpage
\textbf{II. PHẦN TỰ LUẬN}
\begin{ex}
	(1,0 điểm). Cho $a > 0$ và $a \ne 1$. Rút gọn biểu thức
	$P = \left( \dfrac{1}{\sqrt{a}+1} - \dfrac{1}{a+\sqrt{a}} \right) : \dfrac{\sqrt{a}-1}{\sqrt{a}+1}$
	\loigiai{
		Với $a > 0$ và $a \ne 1$, ta có:
		\begin{align*}
			P &= \left( \dfrac{1}{\sqrt{a}+1} - \dfrac{1}{\sqrt{a}(\sqrt{a}+1)} \right) : \dfrac{\sqrt{a}-1}{\sqrt{a}+1} \\
			&= \left( \dfrac{\sqrt{a}}{\sqrt{a}(\sqrt{a}+1)} - \dfrac{1}{\sqrt{a}(\sqrt{a}+1)} \right) \cdot \dfrac{\sqrt{a}+1}{\sqrt{a}-1} \\
			&= \dfrac{\sqrt{a}-1}{\sqrt{a}(\sqrt{a}+1)} \cdot \dfrac{\sqrt{a}+1}{\sqrt{a}-1} \\
			&= \dfrac{1}{\sqrt{a}}
		\end{align*}
		Vậy $P = \dfrac{1}{\sqrt{a}}$.
	}
\end{ex}

\begin{ex}
	(1,0 điểm). Biết phương trình $x^2 - 3x - 5 = 0$ có hai nghiệm $x_1, x_2$. Không giải phương trình, tính giá trị của biểu thức $T = x_1^2 + x_2^2 + x_1 + x_2$.
	\loigiai{
		Phương trình $x^2 - 3x - 5 = 0$ có $a=1, b=-3, c=-5$.
		Vì $ac = 1 \cdot (-5) = -5 < 0$, phương trình luôn có hai nghiệm phân biệt $x_1, x_2$.
		Theo hệ thức Vi-ét, ta có:
		\begin{itemize}
			\item $S = x_1 + x_2 = -\dfrac{b}{a} = -\dfrac{-3}{1} = 3$
			\item $P = x_1 x_2 = \dfrac{c}{a} = \dfrac{-5}{1} = -5$
		\end{itemize}
		Ta có:
		\begin{align*}
			T &= x_1^2 + x_2^2 + x_1 + x_2 \\
			&= (x_1 + x_2)^2 - 2x_1 x_2 + (x_1 + x_2) \\
			&= S^2 - 2P + S \\
			&= (3)^2 - 2(-5) + 3 \\
			&= 9 + 10 + 3 \\
			&= 22
		\end{align*}
		Vậy $T = 22$.
	}
\end{ex}

\begin{ex}
	(1,0 điểm). Giải hệ phương trình $\begin{cases} 3x - y = 5 \\ 4x + y = 9 \end{cases}$
	\loigiai{
		Ta có:
		$$ \begin{cases} 3x - y = 5 \quad (1) \\ 4x + y = 9 \quad (2) \end{cases} $$
		Cộng vế theo vế hai phương trình (1) và (2) của hệ, ta được:
		\begin{align*}
			(3x - y) + (4x + y) &= 5 + 9 \\
			7x &= 14 \\
			x &= \dfrac{14}{7} \\
			x &= 2
		\end{align*}
		Thay $x=2$ vào phương trình (2), ta được:
		\begin{align*}
			4(2) + y &= 9 \\
			8 + y &= 9 \\
			y &= 9 - 8 \\
			y &= 1
		\end{align*}
		Vậy nghiệm của hệ phương trình là $(x; y) = (2; 1)$.
	}
\end{ex}
\begin{ex}
	Một hộp chứa 40 thẻ cùng loại, mỗi thẻ được đánh một số lần lượt từ 1 đến 40. Rút ngẫu nhiên một thẻ trong hộp đó (mỗi thẻ có xác suất như nhau để được chọn). Tính xác suất của biến cố: “Số ghi trên thẻ rút được là bội số của 9”.
	\loigiai{
		Gọi A là biến cố: “Số ghi trên thẻ rút được là bội số của 9”.
		Không gian mẫu là tập hợp các số từ 1 đến 40, do đó số phần tử của không gian mẫu là $n(\Omega) = 40$.
		
		Ta liệt kê các phần tử của biến cố A, tức là các số là bội của 9 trong khoảng từ 1 đến 40:
		$A = \{9; 18; 27; 36\}$
		Số phần tử của biến cố A là $n(A) = 4$.
		
		Xác suất của biến cố A là:
		$$P(A) = \frac{n(A)}{n(\Omega)} = \frac{4}{40} = \frac{1}{10}$$
		Vậy, xác suất để số ghi trên thẻ rút được là bội số của 9 là $\dfrac{1}{10}$.
	}
\end{ex}
\begin{ex}
	Một tổ công nhân theo kế hoạch phải làm 200 sản phẩm trong một thời gian nhất định. Nhưng khi làm thực tế năng suất của tổ đã vượt năng suất dự định là 5 sản phẩm mỗi ngày. Do đó, tổ đã hoàn thành công việc sớm hơn dự định 2 ngày. Hỏi thực tế mỗi ngày tổ đã làm được bao nhiêu sản phẩm?
	\loigiai{
		Gọi số sản phẩm tổ công nhân dự định làm mỗi ngày là $x$ (sản phẩm, $x > 0$).
		Thời gian dự định hoàn thành 200 sản phẩm là $\dfrac{200}{x}$ (ngày).
		
		Thực tế, mỗi ngày tổ làm được $x + 5$ (sản phẩm).
		Thời gian thực tế hoàn thành 200 sản phẩm là $\dfrac{200}{x+5}$ (ngày).
		
		Vì tổ hoàn thành công việc sớm hơn dự định 2 ngày, ta có phương trình:
		$$ \dfrac{200}{x} - \dfrac{200}{x+5} = 2 $$
		Giải phương trình:
		\begin{align*}
			\frac{200(x+5) - 200x}{x(x+5)} &= 2 \\
			\frac{200x + 1000 - 200x}{x^2+5x} &= 2 \\
			\frac{1000}{x^2+5x} &= 2 \\
			1000 &= 2(x^2+5x) \\
			500 &= x^2+5x \\
			x^2 + 5x - 500 &= 0
		\end{align*}
		Ta có $\Delta = b^2 - 4ac = 5^2 - 4(1)(-500) = 25 + 2000 = 2025$.
		Vì $\Delta > 0$, phương trình có hai nghiệm phân biệt:
		$x_1 = \dfrac{-b + \sqrt{\Delta}}{2a} = \dfrac{-5 + \sqrt{2025}}{2(1)} = \dfrac{-5 + 45}{2} = \dfrac{40}{2} = 20$ (thỏa mãn $x>0$)
		$x_2 = \dfrac{-b - \sqrt{\Delta}}{2a} = \dfrac{-5 - \sqrt{2025}}{2(1)} = \dfrac{-5 - 45}{2} = \dfrac{-50}{2} = -25$ (không thỏa mãn $x>0$)
		
		Vậy, số sản phẩm tổ dự định làm mỗi ngày là 20 sản phẩm.
		Số sản phẩm thực tế tổ làm được mỗi ngày là $x+5 = 20+5 = 25$ (sản phẩm).
		
		Đáp số: Thực tế mỗi ngày tổ đã làm được 25 sản phẩm.
	}
\end{ex}
\begin{ex}
	Cho đường tròn $(O; R)$ và dây cung $AB$ khác đường kính. Gọi $H$ là trung điểm của $AB$, đường thẳng $HO$ cắt đường tròn $(O)$ tại hai điểm $C$ và $D$ ($C$ thuộc cung nhỏ $AB$). Điểm $M$ thuộc đoạn $AH$ ($M$ khác $A$ và $H$), tia $CM$ cắt đường tròn $(O)$ tại điểm thứ hai là $N$.
	\begin{enumerate}
		\item[a)] Chứng minh 4 điểm $D, H, M, N$ nằm trên một đường tròn.
		\item[b)] Gọi $I$ là giao điểm của hai đường thẳng $DN$ và $AB$, $K$ là giao điểm của hai đường thẳng $CI$ và $DM$. Chứng minh $NM$ là đường phân giác trong của góc $\widehat{KNH}$ và $M$ là tâm đường tròn nội tiếp tam giác $HKN$.
	\end{enumerate}
\end{ex}

\begin{ex}
	Cho $x, y, z$ là các số thực dương và thỏa mãn điều kiện $x+y+z = xyz$. Tìm giá trị lớn nhất của biểu thức $P = \dfrac{1}{\sqrt{1+x^2}} + \dfrac{1}{\sqrt{1+y^2}} + \dfrac{1}{\sqrt{1+z^2}}$.
\end{ex}

\setcounterref{ex}{0}
